\chapter{Revisión de la literatura}\label{cap:literatura}

Este capítulo revisa, con fuentes verificadas, lo que se sabe sobre cada parte del
problema: CanSat con autogiro, autorrotación vertical de rotores pequeños, sámaras,
aerodinámica de placas curvadas a bajo número de Reynolds, desaceleradores rotativos,
anillo de vórtices, teoría de elemento de pala con correcciones de alta inducción y
paracaídas pequeños. Para cada trabajo se indica qué aporta a este proyecto: un dato, un
método o una advertencia. El capítulo~\ref{cap:arte} ya presentó los antecedentes
básicos; aquí se amplían con 59 referencias nuevas y se contrastan con números del diseño.

\begin{resultado}[Hallazgos principales de la revisión]
\begin{itemize}[leftmargin=*]
  \item \textbf{No hay datos publicados de CanSat con autogiro que midan a la vez masa,
        diámetro, velocidad y rpm.} Los seis trabajos encontrados describen el diseño y,
        a lo sumo, una caída corta. La tabla~\ref{tab:literatura-cansat} queda con muchas
        celdas vacías. Este proyecto puede llenar ese vacío.
  \item \textbf{La placa curvada fina es la elección correcta.} Laitone (1997) midió que
        por debajo de $Re=\num{70000}$ la placa con \SI{5}{\percent} de curvatura en arco
        circular supera a un NACA~0012 en sustentación y en
        $L/D$~\cite{literatura-laitone1997}. La mitad exterior de la pala, que da casi todo
        el empuje, trabaja entre $Re\approx\num{36000}$ y \num{72000}. Martiarena y
        colaboradores midieron en un desacelerador de placas curvadas que la curvatura baja
        las rpm pero no cambia el coeficiente de resistencia
        equivalente~\cite{literatura-martiarena2015}.
  \item \textbf{El paso negativo pequeño está respaldado por la naturaleza y por la
        ingeniería.} Las sámaras naturales operan con paso entre \ang{-0.7} y
        \ang{-2.6}~\cite{literatura-jung2024}; el \emph{Autobody} de la Universidad de
        Maryland usó \ang{-10} con acoplamiento paso--batimiento
        negativo~\cite{literatura-brindejonc2007}. El rango recomendado aquí
        (\ang{-2} a \ang{-4}) queda entre ambos. Ojo: el Autobody necesitó además
        acoplamiento $\delta_3$ y preconicidad negativa para pasar solo a autorrotación;
        este diseño no tiene ninguno de los dos.
  \item \textbf{El punto de operación está donde los modelos empíricos menos coinciden.}
        Con $\CR\approx1{,}2$ la curva de Buhl da un factor de inducción $a\approx0{,}61$
        y el polinomio de Leishman da $a\approx0{,}98$
        (figura~\ref{fig:literatura-ct}). Un balance de energía con la potencia de perfil
        de la pala pide $a\approx0{,}72$--$0{,}91$, entre ambas curvas. Cerrado con una u
        otra curva, el mismo rotor baja a \SIrange{5.5}{6.1}{m/s} (Buhl) o a
        \SIrange{6.5}{6.7}{m/s} (Leishman) (tabla~\ref{tab:literatura-energia}). Hay
        que medir.
  \item \textbf{El drogue da el margen, no la aceptación.} Sin drogue el peor caso de
        diseño, \SI{7.8}{m/s}, ya queda dentro de la banda de \Req{C4}
        (\SIrange{2}{8}{m/s}), pero casi en el borde. Para bajar a \SI{5}{m/s} sin drogue
        haría falta $C_{\mathrm{eq}}\approx2{,}5$ con el disco de \SI{0.4}{m}. El mayor
        valor publicado que se encontró es $\approx1{,}8$, y quizá incluye la resistencia
        del cuerpo.
\end{itemize}
\end{resultado}

%=====================================================================
\section{Método de búsqueda y criterios}\label{sec:literatura-metodo}
La búsqueda se hizo el 5 de octubre de 2026 con un buscador web general, en inglés y en
español, por los ocho temas del encargo. Se aceptó una referencia solo si se pudo
confirmar título, autores, año y medio en la ficha del editor, en un repositorio
institucional o en un índice bibliográfico (PubMed, ADS, NTRS, DTIC, IEEE Xplore). El
volumen, las páginas y el DOI se anotan solo cuando aparecieron en esas fichas; si no, se
omiten.

\begin{nota}[Limitación de acceso]
El entorno de trabajo no permitió descargar los textos completos (los dominios de las
editoriales estaban bloqueados). Los datos numéricos que se citan provienen de los
resúmenes y de las fichas públicas. Antes de usarlos para dimensionar conviene leer el
original. Una revisión posterior volvió a buscar cada cifra en las fichas y resúmenes, y
retiró las que no pudo confirmar. Las celdas vacías de las tablas significan ``no
encontrado'', no ``cero''.
\end{nota}

La figura~\ref{fig:literatura-mapa} ubica en el tiempo las referencias de este capítulo y
las del capítulo~\ref{cap:arte}. Se ven dos olas: la teoría clásica de rotores
(1926--1966) y un crecimiento fuerte después de 2005, impulsado por los microvehículos,
la biomecánica de semillas y la energía eólica. Los CanSat con autogiro son todos
posteriores a 2017.

\begin{figure}[H]
\centering
\begin{tikzpicture}[x=0.108cm, y=0.6cm, every node/.style={font=\scriptsize}]
  % ejes
  \foreach \yr in {1920,1940,1960,1980,2000,2020}{
    \draw[cgris!50] ({\yr-1920},0.4) -- ({\yr-1920},8.6);
    \node[below] at ({\yr-1920},0.4) {\yr};}
  \draw[-{Latex}] (0,0.4) -- (108,0.4) node[right] {año};
  % filas
  \foreach \row/\name in {8/{CanSat con autogiro},7/{Autorrotación pequeña},
      6/{Sámaras y monocópteros},5/{Bajo Reynolds},4/{Desaceleradores rotativos},
      3/{Anillo de vórtices},2/{BEMT y alta inducción},1/{Paracaídas}}{
    \node[left, align=right] at (-1,\row) {\name};
    \draw[cgris!30] (0,\row) -- (106,\row);}
  % CanSat
  \foreach \yr in {2017,2019,2019,2022,2022,2023,2025,2025}{\fill[cfijo] ({\yr-1920},8) circle (2.6pt);}
  % autorrotación pequeña
  \foreach \yr in {1979,1983,2005,2006,2007,2014,2014,2015,2022,2022}{\fill[crot] ({\yr-1920},7) circle (2.6pt);}
  % sámaras
  \foreach \yr in {1973,1989,1992,2007,2009,2010,2010,2012,2013,2014,2015,2016,2019,2022,2022,2024,2024,2026}{\fill[ccuerda] ({\yr-1920},6) circle (2.6pt);}
  % bajo Re
  \foreach \yr in {1942,1979,1983,1995,1996,1997,2000,2003,2008,2011}{\fill[cfreno] ({\yr-1920},5) circle (2.6pt);}
  % desaceleradores
  \foreach \yr in {1966,2004,2013,2024,2024,2024}{\fill[ccont] ({\yr-1920},4) circle (2.6pt);}
  % VRS
  \foreach \yr in {1926,1951,1951,1966,2004,2005,2005,2005,2011}{\fill[cfijo!60!black] ({\yr-1920},3) circle (2.6pt);}
  % BEMT
  \foreach \yr in {1926,1974,2005,2005,2005,2006,2014,2015,2024}{\fill[crot!70!black] ({\yr-1920},2) circle (2.6pt);}
  % paracaídas
  \foreach \yr in {1965,1971,1978,1987,1992,2005,2012,2023}{\fill[ccuerda!60!black] ({\yr-1920},1) circle (2.6pt);}
  % etiquetas de hitos
  \node[above=1pt] at (6,3) {Glauert};
  \node[above=1pt] at (31,3) {Drees};
  \node[above=1pt] at (87,7) {Brindejonc};
  \node[above=1pt] at (89,6) {Lentink};
  \node[above=1pt] at (77,5) {Laitone};
  \node[above=1pt] at (85,2) {Buhl};
  \node[above=1pt] at (46,4) {Barzda};
\end{tikzpicture}
\caption{Distribución temporal de las referencias revisadas (este capítulo y
capítulo~\ref{cap:arte}). Cada punto es una publicación. La asignación de cada trabajo
a una fila es aproximada, porque varios tratan más de un tema.}
\label{fig:literatura-mapa}
\end{figure}

%=====================================================================
\section{CanSat con autogiro}\label{sec:literatura-cansat}
La literatura de CanSat con autogiro es escasa, reciente y casi toda descriptiva: ningún
trabajo encontrado publica a la vez masa, diámetro del rotor, velocidad medida y rpm. La
fuente principal de soluciones probadas siguen siendo las guías y los informes de diseño
de la CanSat Competition~\cite{cansat2019,cansat2025}.

\paragraph{Grupo de la PUCP (Perú).} Rivadeneira y colaboradores presentaron
\emph{Run2Sat~I} en la conferencia IEEE C3 2023~\cite{literatura-rivadeneira2023}: un
CanSat con autogiro que protege un huevo durante el vuelo y el aterrizaje, desarrollado
para el Concurso Iberoamericano de Satélites Enlatados 2023. Informan una eficiencia de
reducción de velocidad de \SI{84.21}{\percent}, sin publicar la velocidad final ni las
rpm. En 2025 publicaron la versión de revista con telemetría por AWS
IoT~\cite{literatura-rivadeneira2025}. Las palas se inspiran en sámaras, y dos pares de
palas giran en sentidos opuestos sobre el mismo eje, cada par sobre su rodamiento.
\emph{Aporte}: la solución contrarrotante evita que el cuerpo gire, a costa de duplicar
el cubo.

\paragraph{Villacorta y colaboradores (2025).} Proponen un autogiro de seis palas para
aterrizar CanSat con cargas frágiles, con el objetivo explícito de reducir la fuerza de
impacto~\cite{literatura-villacorta2025}. Lo validaron con caídas libres desde drones y
edificios, con viento casi nulo; el resumen informa que la velocidad bajó, sin dar cifras. \emph{Aporte}: confirma que el criterio de
diseño útil es la energía o la fuerza de impacto, no solo la velocidad media.

\paragraph{Ventura y colaboradores (2022).} Con perfil NACA~6412 y CFD, buscan bajar de
\SI{34.55}{m/s} a \SI{6}{m/s}; la caída libre de prueba fue desde \SI{39.4}{m} y midió
una desaceleración máxima de \SI{4.94}{m/s^2} en \SI{0.37}{s}~\cite{ventura2022}.
\emph{Advertencia}: una caída de \SI{39}{m} no alcanza para estabilizar la velocidad de
un rotor que arranca a más de \SI{13}{m/s}; el plan de ensayos de este proyecto pide
caídas de 80 a \SI{100}{m} (capítulo~\ref{cap:ensayos}).

\paragraph{Ramadhan y colaboradores (2019) y Shukla y colaboradores (2022).} Ramadhan
describe un CanSat con estructura de PLA con fibra de carbono para la misión 2019: el
contenedor baja con paracaídas a unos \SI{20}{m/s} y a \SI{450}{m} libera la carga con
autogiro, que debe bajar a 10--\SI{15}{m/s}~\cite{ramadhan2019}. Las fuentes accesibles
no dan la masa ni la velocidad medida. Shukla diseña para la misión 2021 dos cargas
autorrotantes con mecanismo de semilla de arce~\cite{literatura-shukla2022}. Kizilkaya y
colaboradores~\cite{literatura-kizilkaya2017} tratan el control de descenso en la misión
2016 (planeador), útil como referencia de método pero no de autogiro.

\paragraph{Informes de equipos.} Los CDR de la misión 2019 son públicos. El del equipo
SwifTurk~\cite{literatura-swifturk2019} describe ocho palas de \SI{270}{mm}. Una rueda
impresa movida por un solo servo abre la tapa inferior y además libera las palas. Un
sensor Hall (Allegro A1104) mide las rpm. \emph{Aporte}: un solo actuador para dos
funciones, y el sensor Hall como solución habitual para medir el giro, igual que aquí.

\begin{table}[H]
\centering\footnotesize
\caption{CanSat con autogiro: datos publicados. ``---'' = no encontrado en las fuentes
accesibles. Las velocidades de 2019 son el requisito de la misión, no una medición.}
\label{tab:literatura-cansat}
\begin{tabularx}{\textwidth}{@{}l c Y c c c c@{}}
\toprule
\textbf{Trabajo} & \textbf{Año} & \textbf{Configuración} & \textbf{Palas} &
\textbf{Masa} & \textbf{$V$ objetivo} & \textbf{rpm}\\\midrule
Ramadhan~\cite{ramadhan2019} & 2019 & Autogiro, misión CanSat 2019 & --- & --- & 10--\SI{15}{m/s} & ---\\
SwifTurk~\cite{literatura-swifturk2019} & 2019 & Palas de \SI{270}{mm}, un servo & 8 & --- & 10--\SI{15}{m/s} & ---\\
Ventura~\cite{ventura2022} & 2022 & NACA~6412, CFD y caída de \SI{39}{m} & --- & --- & \SI{6}{m/s} & ---\\
Shukla~\cite{literatura-shukla2022} & 2022 & Dos cargas tipo semilla de arce & 1 c/u & --- & --- & ---\\
Run2Sat~\cite{literatura-rivadeneira2023,literatura-rivadeneira2025} & 2023--25 & Dos pares contrarrotantes & 2+2 & --- & --- & ---\\
Villacorta~\cite{literatura-villacorta2025} & 2025 & Seis palas; caídas desde drone y edificio & 6 & --- & --- & ---\\\midrule
\textbf{Este diseño} & 2026 & Placa curvada \SI{7}{\percent}, bisagra, drogue atado & 4 & \SI{500}{g} & \SI{6.4}{m/s}\,$^{a}$ & \num{1150}\,$^{a}$\\
\bottomrule
\end{tabularx}

\smallskip\raggedright\footnotesize $^{a}$ Estimación (capítulo~\ref{cap:dim}), no medida.
\end{table}

%=====================================================================
\section{Autorrotación vertical de rotores pequeños}\label{sec:literatura-autorrot}
La autorrotación de cuerpos pequeños tiene una teoría madura y pocos datos de rotores de
menos de \SI{0.5}{m}. Los trabajos más útiles para este proyecto son el \emph{Autobody}
de Maryland, que es casi el mismo problema a mayor escala, y la serie del pararrotor de
la Universidad Politécnica de Madrid, que usa precisamente palas de placa curvada.

\paragraph{Fundamentos.} Lugt (1983) revisa la autorrotación de placas y cuerpos libres y
separa los mecanismos de par impulsor~\cite{literatura-lugt1983}. Iversen (1979) correlacionó
la relación de velocidad de punta de placas planas autorrotantes con su espesor relativo
y su relación de aspecto, y estudió el efecto del momento de
inercia~\cite{literatura-iversen1979}. Seter y Rosen (2014) dan un modelo simplificado
de autorrotación axial de desaceleradores rotativos simples, con los parámetros que
gobiernan el fenómeno, y lo contrastan con un modelo multicuerpo y con
ensayos~\cite{literatura-seter2014}. \emph{Aporte}: el método de Seter y Rosen es el
siguiente paso natural después del modelo de elemento de pala del
capítulo~\ref{cap:teoria}.

\paragraph{El Autobody (Maryland, 2007).} Brindejonc, Sirohi y Chopra diseñaron un
sistema de entrega de cargas por autorrotación con rotor de cuatro palas de
\SI{1.22}{m} de diámetro (\SI{4}{ft}) y masa total de \SI{2.27}{kg}, con requisito de
velocidad máxima de \SI{4.57}{m/s}~\cite{literatura-brindejonc2007}. Lograron la
transición pasiva a autorrotación con un cubo de acoplamiento paso--batimiento negativo
($\delta_3=\ang{-41}$), paso colectivo fijo de \ang{-10} y preconicidad de \ang{-4}; en
caída desde un globo midieron \SI{4.11}{m/s}. Jain, Sirohi y Cameron (2022) extendieron
la idea a un decelerador de despliegue pasivo~\cite{literatura-jain2022}.
\emph{Aporte}: es la validación más cercana de que un rotor libre de paso fijo negativo
arranca solo y se estabiliza sin control activo.

\begin{nota}[Diferencia clave con el Autobody]
El Autobody no confió la transición solo al paso negativo: usó además acoplamiento
paso--batimiento negativo y preconicidad negativa, elegidos con ensayos en túnel
justamente para lograr la transición pasiva~\cite{literatura-brindejonc2007}. Este
diseño tiene bisagra simple, sin $\delta_3$ y sin preconicidad. El par de arranque
positivo con $\theta<0$ (ecuación~\eqref{eq:Q0}) es condición necesaria, pero el
antecedente indica que el paso de arranque a régimen puede requerir más. El ensayo de
arranque en banco (E2, capítulo~\ref{cap:ensayos}) debe verificarlo.
\end{nota}

\paragraph{El pararrotor (UPM).} Nadal Mora, Sanz-Andrés y Cuerva estudiaron en túnel un
desacelerador de alas autorrotantes de baja relación de aspecto y bajo
Reynolds~\cite{literatura-nadal2005} y propusieron un modelo semiempírico de su
comportamiento~\cite{literatura-nadal2006}. Piechocki y colaboradores definieron índices
de estabilidad dinámica cuando el centro de masa está fuera del plano de
palas~\cite{literatura-piechocki2014}, que es la situación de este diseño (cuerpo
colgando bajo el rotor). Martiarena, Nadal Mora y Piechocki midieron en túnel (paso de
\ang{2} a \ang{8}, viento de 4 a \SI{11}{m/s}) el efecto de la \textbf{curvatura de
placa} y de la relación de aspecto~\cite{literatura-martiarena2015}. La relación entre
velocidad de caída y velocidad de giro \textbf{aumenta con la curvatura}, y el
coeficiente de resistencia equivalente \textbf{no cambia} con la curvatura ni con la
relación de aspecto. \emph{Aporte}: es el antecedente experimental más directo para la
placa curvada. Respalda el argumento de la ecuación~\eqref{eq:omega} (más curvatura,
menos rpm a igual caída) y advierte que la curvatura no baja la velocidad de descenso:
su beneficio es mecánico (menos rpm y menos fuerza centrífuga en bisagras y
portapalas), no de frenado. Su
pararrotor tiene alas de baja relación de aspecto, así que los valores no se trasladan
directamente.

\paragraph{Rotores pequeños en descenso axial.} Veismann, Yos y Gharib (2022) hicieron un
estudio paramétrico de rotores pequeños en descenso axial variando paso colectivo,
cuerda, conicidad en planta, número de palas y forma de
punta~\cite{literatura-veismann2022}. Las palas de mayor alargamiento y mayor
coeficiente de carga sufrieron menos pérdida de empuje y menos vibración en el anillo de
vórtices. \emph{Aporte}: es la base de datos moderna más cercana para contrastar el
efecto de $B=4$ y $\sigma=0{,}286$. Ojo: las palas de este diseño tienen alargamiento
bajo ($156/45\approx3{,}5$).

\subsection{Comparación cuantitativa: carga de disco y velocidad}
\label{sec:literatura-carga}
Con la definición de la ecuación~\eqref{eq:CR}, cada decelerador queda caracterizado por
su carga de disco $DL=T/A_R$ y su coeficiente equivalente:
\begin{equation}
  V=\sqrt{\frac{2\,DL}{\rho\,C_{\mathrm{eq}}}},\qquad
  C_{\mathrm{eq}}=\frac{2\,DL}{\rho V^2}.
  \label{eq:literatura-ceq}
\end{equation}
La tabla~\ref{tab:literatura-ceq} aplica esta relación a los únicos datos completos
encontrados (el Autobody) y a este diseño. La figura~\ref{fig:literatura-dl} los ubica
sobre las rectas de $C_{\mathrm{eq}}$ constante.

\begin{table}[H]
\centering\small
\caption{Coeficiente equivalente con área del disco, ecuación~\eqref{eq:literatura-ceq},
$\rho=\SI{1.225}{kg/m^3}$. Script: \texttt{analysis/literatura\_carga\_disco.py}.}
\label{tab:literatura-ceq}
\begin{tabularx}{\textwidth}{@{}Y c c c@{}}
\toprule
\textbf{Caso} & \textbf{$DL$ [\si{N/m^2}]} & \textbf{$V$ [\si{m/s}]} & \textbf{$C_{\mathrm{eq}}$}\\\midrule
Autobody, requisito~\cite{literatura-brindejonc2007} & 19,1 & 4,57 & 1,49\\
Autobody, medido en caída~\cite{literatura-brindejonc2007} & 19,1 & 4,11 & 1,84\\
Este diseño, rotor con drogue ($T=\SI{3.78}{N}$, estimado) & 30,1 & 6,4 & 1,20\\
Este diseño, peor caso sin drogue (estimado) & 39,0 & 7,8 & 1,05\\
Este diseño, sin drogue y $V=\SI{5}{m/s}$ (necesario) & 39,0 & 5,0 & 2,55\\
\bottomrule
\end{tabularx}
\end{table}

\begin{figure}[H]
\centering
\begin{tikzpicture}
\begin{loglogaxis}[width=0.86\textwidth, height=7.2cm,
  xlabel={Carga de disco $DL$ [\si{N/m^2}]}, ylabel={Velocidad de descenso $V$ [\si{m/s}]},
  xmin=5, xmax=120, ymin=1.5, ymax=15,
  log basis x=10, log basis y=10,
  xtick={5,10,20,50,100}, xticklabels={5,10,20,50,100},
  ytick={2,3,5,7,10,15}, yticklabels={2,3,5,7,10,15},
  grid=both, grid style={cgris!25},
  legend columns=2, legend style={font=\scriptsize, cells={anchor=west},
    at={(0.5,-0.2)}, anchor=north}]
  \fill[ccuerda!12] (axis cs:5,2) rectangle (axis cs:120,8);
  \node[font=\scriptsize, ccuerda!70!black, anchor=south east] at (axis cs:118,2.05) {banda \Req{C4}: \SIrange{2}{8}{m/s}};
  \addplot[cgris, dashed, thick] table[x=DL, y=C08] {analysis/data/literatura_curvas.dat};
  \addlegendentry{$C_{\mathrm{eq}}=0{,}8$ (paracaídas plano)}
  \addplot[cfijo, thick] table[x=DL, y=C10] {analysis/data/literatura_curvas.dat};
  \addlegendentry{$C_{\mathrm{eq}}=1{,}0$}
  \addplot[crot, thick] table[x=DL, y=C12] {analysis/data/literatura_curvas.dat};
  \addlegendentry{$C_{\mathrm{eq}}=1{,}2$}
  \addplot[cfreno, thick, densely dotted] table[x=DL, y=C18] {analysis/data/literatura_curvas.dat};
  \addlegendentry{$C_{\mathrm{eq}}=1{,}8$}
  \addplot[only marks, mark=square*, mark size=2.6pt, ccont] coordinates {(19.075,4.11)};
  \addlegendentry{Autobody medido}
  \addplot[only marks, mark=square, mark size=2.6pt, ccont] coordinates {(19.075,4.57)};
  \addlegendentry{Autobody requisito}
  \addplot[only marks, mark=*, mark size=2.8pt, crot] coordinates {(30.07,6.4)};
  \addlegendentry{Este diseño, con drogue (est.)}
  \addplot[only marks, mark=o, mark size=2.8pt, cfijo] coordinates {(39.02,7.8)};
  \addlegendentry{Este diseño, peor caso (est.)}
  \addplot[only marks, mark=x, mark size=3.5pt, cfreno, thick] coordinates {(39.02,5.0)};
  \addlegendentry{$V=\SI{5}{m/s}$ sin drogue}
\end{loglogaxis}
\end{tikzpicture}
\caption{Velocidad de descenso frente a carga de disco. Las rectas son
$C_{\mathrm{eq}}$ constante. Solo se grafican casos con datos publicados completos.}
\label{fig:literatura-dl}
\end{figure}

La lectura es doble. Primero, el diseño propio está en la zona de $\CR\approx1{,}0$--$1{,}2$
que la teoría considera normal para un rotor en autorrotación (capítulo~\ref{cap:teoria}).
Segundo, el Autobody midió un coeficiente equivalente de \num{1.84}, por encima de
$\CR\approx1{,}2$--$1{,}4$, el rango que da el polinomio empírico para el rotor solo en
autorrotación ($V/\vh\approx1{,}7$--$1{,}8$)~\cite{leishman2006,johnson1980}. La diferencia puede deberse a la resistencia del
cuerpo (incluida en la medición, no en el área del disco), a la incertidumbre de la caída
o a la alta solidez del rotor; con los datos accesibles no se puede separar.

\begin{nota}[Qué significa para este diseño]
El valor del Autobody no autoriza a suponer $\CR>1{,}2$ en el dimensionamiento. Sí indica
que el ensayo E3 (capítulo~\ref{cap:ensayos}) puede encontrar un $\CR$ algo mayor que el
supuesto, lo que bajaría $\Vr$. En sentido contrario, alcanzar \SI{5}{m/s} sin el drogue
exigiría $C_{\mathrm{eq}}\approx2{,}55$, valor que ningún trabajo revisado alcanza. Sin
drogue el peor caso (\SI{7.8}{m/s}) cumple \Req{C4} con solo \SI{0.2}{m/s} de margen: el
drogue atado no es redundante, es el que da margen.
\end{nota}

%=====================================================================
\section{Sámaras y robots inspirados en sámaras}\label{sec:literatura-samaras}
Las sámaras demuestran que un rotor de paso fijo, muy fino y algo curvado arranca solo,
se estabiliza solo y es poco sensible a perturbaciones. Las tres lecciones que se
trasladan al diseño son el paso negativo pequeño, la conicidad moderada y la robustez
frente a cambios de masa.

\begin{table}[H]
\centering\footnotesize
\caption{Sámaras y monocópteros: qué aporta cada trabajo a este proyecto.}
\label{tab:literatura-samaras}
\begin{tabularx}{\textwidth}{@{}>{\raggedright\arraybackslash}p{3.1cm} >{\raggedright\arraybackslash}Y >{\raggedright\arraybackslash}Y@{}}
\toprule
\textbf{Trabajo} & \textbf{Resultado} & \textbf{Uso en este proyecto}\\\midrule
Jung y Rezgui 2024~\cite{literatura-jung2024} & Sámaras naturales operan con paso de \ang{-0.7} a \ang{-2.6}; conicidad de \ang{5} a \ang{15} aumenta el empuje vertical hasta \SI{9.6}{\percent}. & Respalda el paso de \ang{-2} a \ang{-4} y la conicidad de \ang{4.4} (dato).\\
Jung y Rezgui 2016~\cite{literatura-jung2016} & El desempeño vertical del rotor tipo sámara empieza a degradarse a escala 8:1; las escaladas autorrotan hasta \ang{80} de ángulo de eje. & Advertencia: no extrapolar coeficientes de semillas a palas de \SI{156}{mm}.\\
Lolli y col.\ 2026~\cite{literatura-lolli2026} & Simulación de 6 grados de libertad de una sámara impresa a escala 0,5, 1 y 2, desde la caída libre hasta la autorrotación estable: la altura de caída necesaria para estabilizarse crece con el tamaño. & Advertencia: un rotor mayor que una semilla necesita más caída para estabilizarse; respalda las caídas largas de los ensayos. Método: CFD de cuerpo libre.\\
Schaeffer y col.\ 2024~\cite{literatura-schaeffer2024} & Cambios de masa mayores al \SI{100}{\percent} cambian la velocidad menos de \SI{15}{\percent}; un modelo de arrastre clásico une las especies. & Contraste: la sámara se autoajusta; en este rotor de paso fijo $V\propto\sqrt{W}$, así que cada gramo de lastre sí sube $\Vr$.\\
El Makdah y col.\ 2022~\cite{literatura-elmakdah2022} & Ante ráfagas verticales, la relación de velocidad de punta sube como máximo \SI{11}{\percent}; el vórtice de borde de ataque sigue adherido. & Advertencia moderada: el viento lateral no detiene el rotor, pero cambia rpm.\\
Lentink y col.\ 2009~\cite{lentink2009}; Lee y col.\ 2014~\cite{literatura-lee2014}; Salcedo y col.\ 2013~\cite{literatura-salcedo2013} & Vórtice de borde de ataque estable en la parte interior de la pala; PIV estereoscópico del flujo 3D. & Explica la sustentación alta de la zona interna de baja velocidad.\\
Varshney y col.\ 2012~\cite{literatura-varshney2012} & Cinemática de la transición de caída a hélice. & Método: modelo de arranque desde reposo.\\
Ulrich y col.\ 2010~\cite{literatura-ulrich2010a,literatura-ulrich2010b}; Kellas 2007~\cite{literatura-kellas2007} & Sámaras robóticas: control de cabeceo y ascenso; sámara guiada para entrega de cargas. & Método: dinámica de vuelo de cuerpo rotante; opción de guiado futuro.\\
Matič y col.\ 2015~\cite{literatura-matic2015} & Modelo de monocóptero con BEMT no estacionaria. & Método directamente aplicable a la dinámica de arranque.\\
Bai y col.\ 2019, 2022~\cite{literatura-bai2019,literatura-bai2022} & Robot de ala giratoria propulsado, inspirado en sámaras, con estabilidad pasiva de actitud. & Referencia de estabilidad pasiva de un cuerpo que gira entero (aquí solo gira el rotor).\\
\bottomrule
\end{tabularx}
\end{table}

\paragraph{Escala.} Una sámara de arce tiene envergadura del orden de centímetros y baja
a velocidades de \SIrange{0.5}{1.5}{m/s} (capítulo~\ref{cap:arte}). El rotor de este
proyecto tiene $R=\SI{0.2}{m}$ y baja de 4 a 13 veces más rápido. Jung y Rezgui
encontraron degradación del desempeño ya a escala 8:1~\cite{literatura-jung2016}, de
modo que las sámaras sirven como guía cualitativa (paso, conicidad, estabilidad), no como
fuente de coeficientes.

%=====================================================================
\section{Aerodinámica de placas curvadas a bajo número de Reynolds}
\label{sec:literatura-lowre}
La evidencia experimental favorece la placa fina curvada frente a perfiles gruesos en el
rango $Re=\num{20000}$--\num{100000}. Ese es justamente el rango del rotor:
$Re\approx\num{36000}$ a media pala y $\approx\num{72000}$ en la punta.

\begin{itemize}
  \item \textbf{Laitone (1997)} ensayó alas rectangulares a números de Reynolds de hasta
        solo $Re=\num{20000}$, en un túnel de baja turbulencia. Para
        todo $Re<\num{70000}$ la mejor sección fue la placa fina con \SI{5}{\percent} de
        curvatura en arco circular: mayor $L/D$ y mayor $C_l$ a todo ángulo. Las placas
        resultaron poco sensibles al número de Reynolds y a la turbulencia del túnel,
        mientras que el NACA~0012 se degradó mucho por debajo de
        $Re=\num{50000}$~\cite{literatura-laitone1997}. \emph{Aporte}: es el dato más
        fuerte a favor de la placa de carbono curvada. Laitone no ensayó el NACA~6412,
        pero la degradación del NACA~0012 sugiere evitar perfiles gruesos a este tamaño.
  \item \textbf{Okamoto y Azuma (2011)} midieron en túnel el efecto de la forma en planta
        a $Re\approx10^4$~\cite{literatura-okamoto2011}, complementando las mediciones de alas de
        libélula~\cite{okamoto1996} y las de Pelletier y Mueller~\cite{pelletier2000}.
  \item \textbf{Bruining (1979)} midió una placa curvada en arco circular, con
        \SI{10}{\percent} de curvatura, a $Re=\num{60000}$ y \num{100000} entre \ang{-10}
        y \ang{90} de ataque~\cite{literatura-bruining1979}. \emph{Aporte}: es de los
        pocos conjuntos que cubre ángulos grandes, que son los que ve la pala al
        arrancar (capítulo~\ref{cap:teoria}, par de arranque).
  \item \textbf{Shyy y colaboradores (2008)} sintetizan la física de transición, burbuja
        de separación laminar y efectos no estacionarios en voladores de bajo
        Reynolds~\cite{literatura-shyy2008}, junto con las revisiones de
        Lissaman~\cite{lissaman1983} y Mueller y DeLaurier~\cite{mueller2003}.
\end{itemize}

\begin{nota}[Advertencia sobre las polares]
Laitone midió con \SI{5}{\percent} de curvatura y Bruining con \SI{10}{\percent}; el
diseño usa \SI{7}{\percent}, entre ambos. Las polares de la tabla~\ref{tab:polares} son
interpolaciones y cálculos, no mediciones de esta placa. Además, todos esos datos son de
alas que no giran: la rotación y el vórtice de borde de ataque (sección
\ref{sec:literatura-samaras}) pueden aumentar el $C_l$ de la zona interna.
\end{nota}

%=====================================================================
\section{Desaceleradores rotativos para entrada atmosférica y recuperación}
\label{sec:literatura-decel}
Los rotores como freno de reentrada o de recuperación se estudian desde los años 60 y
reaparecen hoy para cohetes reutilizables y sondas planetarias. Su mensaje para este
proyecto es que el rotor libre funciona y es liviano, pero la etapa crítica es el
despliegue y el arranque.

{\footnotesize
\begin{longtable}{@{}>{\raggedright\arraybackslash}p{2.6cm} >{\raggedright\arraybackslash}p{\dimexpr(\textwidth-2.6cm-4\tabcolsep)/2\relax} >{\raggedright\arraybackslash}p{\dimexpr(\textwidth-2.6cm-4\tabcolsep)/2\relax}@{}}
\caption{Desaceleradores rotativos: resumen y aporte.}\label{tab:literatura-decel}\\
\toprule
\textbf{Trabajo} & \textbf{Resultado} & \textbf{Uso en este proyecto}\\\midrule
\endfirsthead
\multicolumn{3}{@{}l}{\emph{Tabla~\thetable\ (continuación)}}\\\toprule
\textbf{Trabajo} & \textbf{Resultado} & \textbf{Uso en este proyecto}\\\midrule
\endhead
\bottomrule
\endlastfoot
Barzda 1966~\cite{literatura-barzda1966} & Desarrollo del \emph{Rotochute} de Kaman: rotores de 0,3 a \SI{7.3}{m} de diámetro, cargas de 2,7 a \SI{408}{kg}, descensos de hasta \SI{6}{m/s}; despliegues hasta Mach~1,2. & Rotor plegado que se abre en vuelo, ensayado incluso con \SI{0.3}{m}.\\
Young y col.\ 2004~\cite{literatura-young2004} & Rotores desplegados en la parte trasera de una sonda, en autorrotación o como turbina, para el descenso en Venus. & Muestra que el rotor puede además generar energía (opción futura de telemetría).\\
Diaz-Silva y col.\ 2013~\cite{literatura-diazsilva2013} & Revisión histórica y simulación; el sistema por rotor pesaría menos del \SI{15}{\percent} del peso de aterrizaje. & Aquí las cuatro palas con portapalas (\SI{\approx30}{g}) son el \SI{6}{\percent} de la masa, sin el cubo.\\
Bergmann y col.\ 2024~\cite{literatura-bergmann2024} & \emph{Daedalus~2} en REXUS~29: dos sondas expulsadas a unos \SI{80}{km}, con rotor pasivo y control de paso. Los autores informan que no lograron el aterrizaje en autorrotación. & Advertencia: aun con control de paso no lograron aterrizar; medir rpm en vuelo es clave para diagnosticar.\\
Marques y col.\ 2024~\cite{literatura-marques2024} & Rotor de recuperación de un cohete sonda: el paso es el parámetro más influyente; más palas o más radio pueden empeorar por peso. & Respalda elegir el paso con cuidado; método de trayectoria.\\
Patnala y col.\ 2024~\cite{literatura-patnala2024} & Descenso axial y helicoidal de un rotor no propulsado con optimización de trayectoria. & Método: puntos fijos de la autorrotación axial.\\
Jain y col.\ 2022~\cite{literatura-jain2022} & Despliegue pasivo de palas y pruebas. & Despliegue sin actuadores de palas plegadas.\\
\end{longtable}
}

%=====================================================================
\section{Anillo de vórtices y estado de estela turbulenta}\label{sec:literatura-vrs}
Un rotor libre que se suelta a \SI{13.4}{m/s} y termina a unos \SI{6}{m/s} pasa del
estado de molino al de estela turbulenta y se estabiliza cerca del borde del anillo de
vórtices, sin entrar en él. La literatura coincide en que en esa zona la teoría de
cantidad de movimiento no vale y en que, dentro del anillo de vórtices, el empuje
fluctúa con fuerza.

\begin{itemize}
  \item \textbf{Glauert (1926)} analizó los estados de freno de molino y anillo de
        vórtices de una hélice y propuso la curva empírica que todavía se usa como
        base~\cite{literatura-glauert1926b}.
  \item \textbf{Drees y Hendal (1951)} visualizaron con humo el flujo alrededor de un
        rotor en túnel y fijaron la imagen clásica del anillo de
        vórtices~\cite{literatura-drees1951}.
  \item \textbf{Washizu y colaboradores (1966)} midieron en vía móvil un rotor de
        \SI{1.1}{m} ($\sigma=\num{0.0573}$): en el anillo de vórtices el empuje fluctúa
        con violencia y el par casi nada. Las fluctuaciones aparecieron entre
        $V/\vh\approx0{,}23$ y $1{,}25$ en descenso axial~\cite{literatura-washizu1966}.
  \item \textbf{Green, Gillies y Brown (2005)} midieron con PIV el campo alrededor de un
        rotor modelo que desciende en su propia estela~\cite{literatura-green2005}.
  \item \textbf{Stack, Caradonna y Savaş (2005)} observaron en tanque de agua que en el
        anillo de vórtices se desprenden anillos periódicamente, con fluctuaciones de
        empuje de pico a pico de hasta \SI{95}{\percent} del valor medio cada 20 a 50
        vueltas~\cite{literatura-stack2005}.
  \item \textbf{Shetty y Selig (2011)} ensayaron 26 hélices pequeñas (9 a 11~pulgadas)
        a bajo Reynolds en descenso vertical, con relación de avance de $-0{,}8$ a 0, y
        registraron las fluctuaciones de empuje~\cite{literatura-shetty2011}.
        \emph{Aporte}: es el conjunto amplio más cercano en escala que se encontró
        (\SIrange{0.23}{0.28}{m} de diámetro, frente a \SI{0.4}{m} aquí).
\end{itemize}

\paragraph{Consecuencia cuantitativa.} A \num{1150} rpm, 20 a 50 vueltas duran de
\num{1.0} a \SI{2.6}{s}. Si el rotor cayera en anillo de vórtices, la carga del cubo
podría oscilar con ese período. La estimación de diseño, $V_2/\vh=1{,}83$, ubica al rotor en
el estado de estela turbulenta, cerca de la autorrotación ideal. Queda fuera de la zona
de fluctuaciones de Washizu ($V/\vh\le1{,}25$), con un margen de
$0{,}58\,\vh\approx\SI{2}{m/s}$. Esta lectura coincide con los modelos del borde del anillo de vórtices de
Leishman y colaboradores y de Johnson~\cite{leishman2004,johnson2005}. Por eso el
arranque a velocidad alta (\SI{13.4}{m/s}, $V/\vh\approx3{,}8$) es favorable: el rotor
entra desde el estado de molino y no desde el vuelo estacionario. Como $V/\vh=2/\sqrt{\CR}$
para el rotor solo, el margen se reduce si el $\CR$ real es mayor que el supuesto: con
$\CR=1{,}75$ (el extremo de la tabla~\ref{tab:literatura-energia}) queda
$V/\vh\approx1{,}51$, todavía fuera de esa zona. El drogue no cambia esta relación.

%=====================================================================
\section{Elemento de pala y correcciones de alta inducción}\label{sec:literatura-bemt}
La teoría combinada de elemento de pala y cantidad de movimiento (BEMT) de las turbinas
eólicas es la herramienta natural para este rotor, que en descenso es un molino de
viento libre. La literatura eólica aporta tres piezas que la teoría de helicópteros trata
menos: corrección de alta inducción, corrección de punta y un método de solución que
siempre converge.

\paragraph{Corrección de alta inducción.} La teoría de cantidad de movimiento da
$C_T=4a(1-a)$, con un máximo de 1 en $a=0{,}5$; las mediciones siguen subiendo. Glauert
propuso una corrección empírica~\cite{literatura-glauert1926b} y Wilson y Lissaman la
llevaron a la práctica de las turbinas~\cite{literatura-wilson1974}. Buhl (2005)
reemplazó la curva por una parábola continua con la rama de cantidad de movimiento en
$a_c=0{,}4$, incluida la pérdida de punta $F$~\cite{literatura-buhl2005}:
\begin{equation}
  C_T=\begin{cases}
    4aF(1-a), & a\le0{,}4,\\[2pt]
    \dfrac89+\left(4F-\dfrac{40}{9}\right)a+\left(\dfrac{50}{9}-4F\right)a^2, & a>0{,}4.
  \end{cases}
  \label{eq:literatura-buhl}
\end{equation}
CCBlade, el código de Ning~\cite{literatura-ning2014}, usa esta corrección; el
manual teórico de AeroDyn~\cite{literatura-moriarty2005} describe el tratamiento del estado
de estela turbulenta en ese código.

\paragraph{Corrección de punta y solución.} Shen y colaboradores mostraron que las
correcciones de punta derivadas de Prandtl son inconsistentes cerca de la punta y
propusieron una nueva~\cite{literatura-shen2005}. Ning (2014) reformuló las ecuaciones
BEMT como una sola ecuación en el ángulo de flujo $\phi$, con convergencia
garantizada~\cite{literatura-ning2014}; esto es valioso aquí porque con paso negativo y
alta inducción los métodos iterativos en $a$ suelen fallar. Hansen~\cite{literatura-hansen2015}
es el texto de referencia. Liew, Heck y Howland (2024) propusieron un modelo de cantidad
de movimiento unificado que predice el empuje con coeficientes altos sin correcciones
empíricas~\cite{literatura-liew2024}.

\paragraph{Contraste con la curva de helicópteros.} La relación empírica de este informe
(ecuación~\eqref{eq:poly}) viene de rotores de helicóptero~\cite{castles1951,leishman2006}.
Puede pasarse a la notación eólica con $a=\vi/V$ y $C_T=4(\vh/V)^2$. La
figura~\ref{fig:literatura-ct} compara las tres curvas.

\begin{figure}[H]
\centering
\begin{tikzpicture}
\begin{axis}[width=0.84\textwidth, height=7cm,
  xlabel={Factor de inducción axial $a=\vi/V$}, ylabel={$C_T=T/(\tfrac12\rho V^2A_R)$},
  xmin=0, xmax=1.2, ymin=0, ymax=2.0, grid=both, grid style={cgris!25},
  legend columns=2, legend style={font=\scriptsize, cells={anchor=west},
    at={(0.5,-0.2)}, anchor=north}]
  \fill[ccuerda!15] (axis cs:0.717,0) rectangle (axis cs:0.906,2.0);
  \node[font=\scriptsize, ccuerda!60!black, align=center] at (axis cs:0.81,0.2) {balance\\de energía};
  \addplot[cgris, thick, dashed] table[x=a, y=CTmom] {analysis/data/literatura_ct_a.dat};
  \addlegendentry{Cantidad de movimiento $4a(1-a)$}
  \addplot[cfijo, thick] table[x=a, y=CTbuhl] {analysis/data/literatura_ct_a.dat};
  \addlegendentry{Buhl (2005), $F=1$}
  \addplot[cfijo, thin, densely dotted] table[x=a, y=CTbuhl09] {analysis/data/literatura_ct_a.dat};
  \addlegendentry{Buhl (2005), $F=0{,}9$}
  \addplot[crot, thick] table[x=a, y=CT] {analysis/data/literatura_ct_leish.dat};
  \addlegendentry{Polinomio de Leishman, ec.~\eqref{eq:poly}}
  \draw[cfreno, thin] (axis cs:0,1.2) -- (axis cs:1.2,1.2);
  \node[cfreno, font=\scriptsize, anchor=south west] at (axis cs:0.02,1.2) {$\CR=1{,}2$ (diseño)};
  \addplot[only marks, mark=*, cfijo, mark size=2.4pt] coordinates {(0.612,1.2)};
  \addplot[only marks, mark=*, crot, mark size=2.4pt] coordinates {(0.982,1.194)};
  \node[font=\scriptsize, cfijo, anchor=south east] at (axis cs:0.60,1.22) {$a=0{,}61$};
  \node[font=\scriptsize, crot, anchor=north west] at (axis cs:0.99,1.17) {$a=0{,}98$};
\end{axis}
\end{tikzpicture}
\caption{Coeficiente de empuje frente a factor de inducción: dos curvas empíricas de
origen distinto (turbinas y helicópteros) dan factores de inducción muy diferentes para el
$\CR$ de diseño. La franja verde es el rango de $a$ que pide el balance de energía de la
tabla~\ref{tab:literatura-energia} (estimación). Datos:
\texttt{analysis/data/literatura\_ct\_*.dat}.}
\label{fig:literatura-ct}
\end{figure}

\paragraph{Cierre por balance de energía.} Las dos curvas dan el mismo $V/\vh$ para un
$\CR$ dado; difieren en cuánto aire atraviesa el disco. Ese flujo lo fija la potencia de
perfil. En autorrotación el rotor no recibe potencia por el eje, así que la potencia que
el aire entrega al disco, $T\,(V-\vi)$, debe igualar la que disipa la resistencia de
perfil de las palas:
\begin{equation}
  T\,(V-\vi)=P_0\approx B\,\tfrac12\rho\,c\,\bar C_d\,\Omega^3\,\frac{R^4-e^4}{4}
  \quad\Rightarrow\quad a=1-\frac{P_0}{T\,V}.
  \label{eq:literatura-energia}
\end{equation}
Con $\Omega=\SI{120}{rad/s}$ (estimación base), $T=\SI{3.78}{N}$ y $\bar C_d$ entre
\num{0.03} y \num{0.09}, el flujo neto por el disco es de \SIrange{0.6}{1.8}{m/s} y
$a\approx0{,}72$--$0{,}91$. Si además se impone cada curva empírica, con el drogue y el
cuerpo en corriente libre, se obtiene la velocidad de equilibrio de la
tabla~\ref{tab:literatura-energia}.

\begin{table}[H]
\centering\small
\caption{Cierre por balance de energía, ecuación~\eqref{eq:literatura-energia}, con
$\Omega=\SI{120}{rad/s}$ fijo. $w=V-\vi$ y $a$ se calculan en el punto base
($T=\SI{3.78}{N}$, $V=\SI{6.4}{m/s}$); $\phi_{75}$ es el ángulo de flujo en $0{,}75R$.
Las columnas de la derecha resuelven el equilibrio completo con cada curva. Estimación de
orden; script \texttt{analysis/literatura\_carga\_disco.py}.}
\label{tab:literatura-energia}
\begin{tabular}{@{}ccccc cc cc@{}}
\toprule
 & & & & & \multicolumn{2}{c}{\textbf{Leishman}} & \multicolumn{2}{c}{\textbf{Buhl}}\\
\cmidrule(lr){6-7}\cmidrule(l){8-9}
$\bar C_d$ & $P_0$ [\si{W}] & $w$ [\si{m/s}] & $a$ & $\phi_{75}$ [\si{\degree}] & $V$ [\si{m/s}] & $\CR$ & $V$ [\si{m/s}] & $\CR$\\\midrule
0,030 & 2,28 & 0,60 & 0,91 & 1,9 & 6,50 & 1,15 & 5,50 & 1,75\\
0,045 & 3,42 & 0,91 & 0,86 & 2,9 & 6,56 & 1,12 & 5,66 & 1,63\\
0,060 & 4,56 & 1,21 & 0,81 & 3,8 & 6,62 & 1,10 & 5,80 & 1,53\\
0,075 & 5,70 & 1,51 & 0,76 & 4,8 & 6,67 & 1,08 & 5,95 & 1,44\\
0,090 & 6,84 & 1,81 & 0,72 & 5,7 & 6,72 & 1,05 & 6,09 & 1,36\\
\bottomrule
\end{tabular}
\end{table}

\begin{resultado}[Consecuencia para el cálculo de rpm y de velocidad]
Para $\CR=1{,}2$ la curva de Buhl da $a=0{,}61$: el aire atraviesa el disco a
$V(1-a)\approx\SI{2.5}{m/s}$. El polinomio de Leishman da $a=0{,}98$, con flujo neto
casi nulo. El balance de energía con la estimación base de rpm pide un valor intermedio,
$a\approx0{,}72$--$0{,}91$, con $\phi_{75}\approx\ang{2}$--\ang{6}. Con
$\theta=\ang{-3}$ eso da $\alpha=\phi+\theta\approx\ang{-1}$ a \ang{3}, donde una placa
con \SI{7}{\percent} de curvatura puede dar $C_l\approx0{,}9$. La estimación de
\num{1150} rpm (ecuación~\eqref{eq:omega}) es entonces coherente, pero no la elige
ninguna de las dos curvas. Cerrado con Buhl, el mismo rotor bajaría a
\SIrange{5.5}{6.1}{m/s}; con Leishman, a \SIrange{6.5}{6.7}{m/s}. Ambos casos cumplen
\Req{C4}; la diferencia, de casi \SI{1}{m/s}, es la incertidumbre de modelo que hoy
tiene $\Vr$. Las mediciones de rpm (E2) y de $\CR$ (E3) deciden.
\end{resultado}

%=====================================================================
\section{Paracaídas pequeños y drogue atado}\label{sec:literatura-paracaidas}
Para paracaídas pequeños los manuales dan $C_D\approx0{,}75$--$0{,}8$ para el
paracaídas plano circular referido al área nominal~\cite{knacke1992}, y los ensayos
advierten sobre dos fenómenos que la estimación cuasi estática ignora: la
``respiración'' de la campana y la oscilación pendular.

\begin{itemize}
  \item \textbf{Manuales.} Knacke~\cite{knacke1992}, la guía de Ewing, Bixby y
        Knacke~\cite{literatura-ewing1978} y el AGARDograph de
        Cockrell~\cite{literatura-cockrell1987} son la base de coeficientes, cargas de
        apertura y estabilidad.
  \item \textbf{Scher y Young (1971)} midieron en el túnel vertical de Langley el
        $C_D$ de paracaídas planos circulares de \SI{1.07}{m}, sostenidos con marcos de
        alambre en las formas que toman durante el inflado. Los de baja porosidad
        resultaron inestables y los de alta porosidad,
        estables~\cite{literatura-scher1971}. \emph{Aporte}: datos para modelar la
        apertura del drogue, y la advertencia de que la tela poco porosa y sin
        ventilación oscila.
  \item \textbf{Johari y Desabrais (2005)} midieron en túnel de agua que la campana
        ``respira'' con frecuencia adimensional $\approx0{,}55$ y desprende anillos de
        vórtice a esa misma frecuencia; la masa añadida explica hasta el
        \SI{40}{\percent} de las fluctuaciones de resistencia~\cite{literatura-johari2005}.
        La frecuencia está referida a la velocidad libre y al diámetro
        \emph{proyectado} medio de la campana. Si se escala igual, para el drogue de
        \SI{0.25}{m} de diámetro nominal (proyectado de unos \SI{0.17}{m}) eso es
        $f\approx0{,}55\,V/D_p$, es decir \SIrange{14}{21}{Hz} a \SI{6.4}{m/s} y
        \SIrange{29}{43}{Hz} a \SI{13.4}{m/s} (estimación; el extremo bajo usa el
        diámetro nominal). En la etapa 2 el drogue está en la estela del rotor, así que
        la velocidad de referencia es incierta.
  \item \textbf{Mattei y colaboradores (2023)} encontraron en paracaídas a escala que a
        bajo Reynolds la respiración tiene gran amplitud y baja frecuencia, con grandes
        fluctuaciones de resistencia~\cite{literatura-mattei2023}. \emph{Advertencia}: el
        drogue de este diseño es pequeño y lento; sus fluctuaciones relativas pueden ser
        mayores que las de los manuales.
  \item \textbf{Guglieri (2012)} da un modelo multicuerpo simple de paracaídas y carga
        para simular trayectorias~\cite{literatura-guglieri2012}. \emph{Método}: sirve
        para modelar el conjunto drogue--cuerda--rotor--cuerpo de la etapa 2.
\end{itemize}

\begin{nota}[Interacción drogue--rotor]
No se encontró ningún trabajo que estudie un drogue atado por encima de un rotor
autorrotante, con la cuerda pasando por el eje. En descenso el aire relativo sube: pasa
primero por el cuerpo y el rotor y después llega al drogue, \SI{1.5}{m} (unos $3{,}75R$)
más arriba. El drogue, de \SI{0.25}{m}, queda entonces dentro de la estela del rotor de
\SI{0.4}{m}: un flujo frenado, con giro y turbulento. Con $a\approx0{,}6$--$1$
(figura~\ref{fig:literatura-ct}; $0{,}72$--$0{,}91$ por balance de energía) la velocidad de esa estela es mucho menor que $V$ y puede
incluso recircular, de modo que la presión dinámica sobre el drogue, y su aporte al
freno, son inciertos. Por eso el peor caso del diseño (\SI{7.8}{m/s}) supone drogue sin
aporte. Medir $\Vr$ con y sin drogue es la única forma de cerrar esta duda.
\end{nota}

%=====================================================================
\section{Síntesis: decisiones de diseño y fuerza de la evidencia}
\label{sec:literatura-sintesis}
La tabla~\ref{tab:literatura-trazab} vincula cada decisión del diseño con la evidencia
publicada y califica su fuerza: \textbf{A}~=~medición directa en condiciones parecidas,
\textbf{B}~=~medición en otra escala o configuración, o práctica de diseño documentada
sin medición publicada, \textbf{C}~=~solo teoría o analogía. Ninguna decisión llega a A:
no hay mediciones publicadas de un rotor como este.

{\footnotesize
\begin{longtable}{@{}>{\raggedright\arraybackslash}p{3.3cm} >{\raggedright\arraybackslash}p{\dimexpr\textwidth-3.3cm-1.2cm-4\tabcolsep\relax} >{\centering\arraybackslash}p{1.2cm}@{}}
\caption{Trazabilidad de las decisiones de diseño a la literatura.}\label{tab:literatura-trazab}\\
\toprule
\textbf{Decisión} & \textbf{Evidencia} & \textbf{Fuerza}\\\midrule
\endfirsthead
\multicolumn{3}{@{}l}{\emph{Tabla~\thetable\ (continuación)}}\\\toprule
\textbf{Decisión} & \textbf{Evidencia} & \textbf{Fuerza}\\\midrule
\endhead
\bottomrule
\endlastfoot
Placa de carbono curvada \SI{7}{\percent} & Laitone~\cite{literatura-laitone1997}, Bruining~\cite{literatura-bruining1979}, Pelletier y Mueller~\cite{pelletier2000}; Martiarena~\cite{literatura-martiarena2015}: la curvatura baja las rpm pero no cambia $C_{\mathrm{eq}}$ & B\\
Paso $-2$ a \ang{-4}, sin control & Jung y Rezgui~\cite{literatura-jung2024}, Seter y Rosen~\cite{literatura-seter2014}, Marques~\cite{literatura-marques2024}; el Autobody~\cite{literatura-brindejonc2007} necesitó además $\delta_3$ y preconicidad & B/C\\
$\CR=1{,}2$ para dimensionar & Castles y Gray~\cite{castles1951}, Leishman~\cite{leishman2006}; el Autobody sugiere margen~\cite{literatura-brindejonc2007}; el cierre de energía da 1,05 a 1,75 según la curva & B\\
Bisagra, conicidad \ang{4.4} & Jung y Rezgui (conicidad \ang{5}--\ang{15})~\cite{literatura-jung2024}; teoría~\cite{johnson1980} & B\\
$B=4$, $\sigma=0{,}286$ & Veismann y col.~\cite{literatura-veismann2022} (paramétrico); sin dato directo & C\\
Drogue atado & Ningún antecedente directo; manuales~\cite{knacke1992,literatura-scher1971} & C\\
Arranque a \SI{13.4}{m/s} & Estela turbulenta y anillo de vórtices~\cite{literatura-stack2005,literatura-shetty2011,leishman2004} & B\\
Traba con servo & SwifTurk~\cite{literatura-swifturk2019} (un servo abre la tapa y libera las palas) & B\\
Sensor Hall para rpm & CDR 2019~\cite{literatura-swifturk2019}, guía 2025~\cite{cansat2025} & B\\
\end{longtable}
}

%=====================================================================
\section{Vacíos de conocimiento que este proyecto puede cubrir}\label{sec:literatura-vacios}
La revisión deja cinco vacíos concretos. Cada uno se puede cubrir con ensayos ya
previstos en el capítulo~\ref{cap:ensayos}, sin costo adicional relevante.

\begin{enumerate}
  \item \textbf{Datos completos de un CanSat con autogiro.} Ningún trabajo publica masa,
        diámetro, velocidad estabilizada medida y rpm del mismo vehículo. Publicar esas
        cuatro cifras con su incertidumbre sería, por sí solo, un aporte.
  \item \textbf{$\CR$ medido de un rotor de \SI{0.4}{m} de placa curvada a
        $Re<\num{100000}$.} Los coeficientes disponibles son de helicópteros grandes o de
        sámaras. La tabla~\ref{tab:literatura-ceq} muestra una dispersión de 1,0 a 1,8.
  \item \textbf{Inducción en autorrotación axial a bajo Reynolds.} La
        figura~\ref{fig:literatura-ct} muestra que las curvas empíricas de turbinas y de
        helicópteros discrepan justo en el punto de operación. Medir rpm, $\Vr$ y $\CR$ a la
        vez, junto con la polar de la placa, permite estimar $a$ con la
        ecuación~\eqref{eq:literatura-energia}.
  \item \textbf{Arranque de un rotor plegado liberado a alta velocidad.} Hay estudios de
        arranque de sámaras desde reposo~\cite{literatura-varshney2012} y de despliegue
        pasivo~\cite{literatura-jain2022}, pero no del tiempo de arranque de un rotor
        de cuatro palas con bisagra soltado a \SI{13}{m/s}. La cámara cenit y el sensor
        Hall pueden medirlo.
  \item \textbf{Drogue atado sobre un rotor.} No se encontró ningún estudio de esta
        combinación. Un ensayo con y sin drogue (misma caída) separaría su aporte.
\end{enumerate}

\begin{nota}[Límites de esta revisión]
La revisión es amplia pero no exhaustiva: no incluye literatura en turco, chino o ruso,
donde hay muchos equipos de CanSat y de Teknofest, ni informes de equipos que no estén
indexados. Varias fichas no muestran volumen o páginas; se dejaron sin ellos en vez de
completarlos de memoria. Los números atribuidos a cada trabajo deben confirmarse en el
texto completo antes de una entrega formal.
\end{nota}
